% Backup for DLCoEP, RNN: optional and advanced. Used only by this course, at the end of the deck,
% after the case study (see seminar_dlcoep_rnn_content.tex).
% Relocated 2026-10-26 from the middle of dl_rnn.tex: the author found these three frames too advanced for beginners.
% The text is unchanged. To show them in class, move this file's \input before the case study.

\section[Backup]{Backup: Optional, Advanced}

%%%%%%%%%%%%%%%%%%%%%%%%%%%%%%%%%%%%%%%%%%%%%%%%%%%
\begin{frame}
  \begin{center}
    {\Large Backup: Optional, Advanced}

    {\small Experiments that compare a dense layer with recurrent layers}
  \end{center}
\end{frame}

%%%%%%%%%%%%%%%%%%%%%%%%%%%%%%%%%%%%%%%%%%%%%%%%%%%
\begin{frame}[fragile] \frametitle{Test 1: Order and Length}
\begin{itemize}
\item Task: symbols A and B each appear once among filler symbols. Label: does A come before B? Trained on sequences of length 6
\end{itemize}
\begin{center}
\small
\adjustbox{max width=\linewidth}{%
\begin{tabular}{lrcc}
\hline
Model & Parameters & Test, length 6 & Test, length 10 \\
\hline
Bag of symbols (counts) + dense & 114 & 49.5\% & 49.9\% \\
Dense on the 6 positions & 866 & 100.0\% & cannot run \\
GRU, last state & 1,090 & 100.0\% & 50.1\% \\
\hline
\end{tabular}}
\end{center}
\begin{itemize}
\item The bag cannot tell the two orders apart: a coin flip. The dense net reads order but accepts only length 6
\item The same GRU trained on lengths 6, 8, 10 and 12 scores 100.0\% on lengths 6, 10 and an unseen 14, and 83.1\% on an unseen 20. Running on any length is not enough: it has to be trained on a range of lengths
\end{itemize}
\end{frame}

%%%%%%%%%%%%%%%%%%%%%%%%%%%%%%%%%%%%%%%%%%%%%%%%%%%
\begin{frame}[fragile] \frametitle{Test 2: Move the Spike in Time}
\begin{itemize}
\item Task: is there a spike somewhere in a noisy sequence of 20 steps? Train with the spike only in steps 1 to 10, then test with it in steps 1 to 10 and in steps 11 to 20
\end{itemize}
\begin{center}
\small
\adjustbox{max width=\linewidth}{%
\begin{tabular}{lrcc}
\hline
Model & Parameters & Spike in steps 1 to 10 & Spike in steps 11 to 20 \\
\hline
Dense 20-32-2 & 738 & 100.0\% & 50.2\% \\
GRU, 8 units & 282 & 100.0\% & 97.1\% \\
\hline
\end{tabular}}
\end{center}
\begin{itemize}
\item The dense net's weights for steps 11 to 20 never saw a spike: a coin flip
\item The GRU applies the same weights at every step, so a late spike is still recognised. The small drop (97.1\%) is because a late spike also arrives in a different memory state
\end{itemize}
\end{frame}

%%%%%%%%%%%%%%%%%%%%%%%%%%%%%%%%%%%%%%%%%%%%%%%%%%%
\begin{frame}[fragile] \frametitle{Parameters, and the Three Ideas}
\begin{center}
\small
\adjustbox{max width=\linewidth}{%
\begin{tabular}{rrrrr}
\hline
Steps $T$ & Dense on $T \times 8$ inputs & RNN & GRU & LSTM \\
\hline
10 & 2,592 & 1,344 & 4,032 & 5,376 \\
100 & 25,632 & 1,344 & 4,032 & 5,376 \\
1000 & 256,032 & 1,344 & 4,032 & 5,376 \\
\hline
\end{tabular}}
\end{center}
\begin{itemize}
\item $8$ readings per step and $32$ hidden units. The dense layer grows with $T$ and accepts only one $T$; the recurrent layers do not grow
\item RNN count: $d \times H + H \times H + 2H = 8 \times 32 + 32 \times 32 + 64 = 1{,}344$ (PyTorch keeps two bias vectors)
\item So: any length, order kept in a hidden state, the same weights at every step. That is the recurrent network
\end{itemize}
\end{frame}
